\section{Experiments}

\subsection{Motif-Pair Displacement}

\paragraph*{Task family and shared structure}

We consider binary classification tasks indexed by $\tau=(A,B,g)$, where
$A,B\in\{0,1\}^3$ are distinct ordered binary patterns and $g$ is their
circular displacement in a sequence $x\in\{0,1\}^{16}$. The label is $y=1$
if the unique occurrences of $A$ and $B$ are separated by $g$, and $y=0$
otherwise (Figure~\ref{fig:motif-pair-task}). The classes are balanced, and
negative examples contain both patterns at another admissible displacement
while matching the number of ones. Across tasks, $A$ and $B$ determine
task-specific weights, whereas $g$ selects a member of a shared Toeplitz-like
support family consisting of two circular diagonal bands.

\begin{figure}[htbp]
  \centering
  \includegraphics[width=0.75\linewidth]{motif_pair_A001_B110_G05}
  \caption{Positive and hard-negative examples for a binary-pattern pair
    task. Both sequences contain unique occurrences of $A$ and $B$ and have
    the same number of ones; only their circular displacement determines the
    label.}
  \label{fig:motif-pair-task}
\end{figure}

\paragraph*{Structural-prior extraction}

For each meta-training task $\tau$, we independently sample $N=512$
first-layer masks from the proposal distribution $p_0$. We use a product
Bernoulli distribution with activation probability $\rho=0.375$, chosen so
that
\[
\mathbb{E}_{M\sim p_0}\!\left[\lVert M\rVert_0\right]
=256\rho=96,
\]
matching the number of active connections in the reference and generated
structures. Given a mask $M$, we train from scratch a one-hidden-layer MLP
\[
f_{M,w}(x)
=
W^{(2)}\sigma\!\left((W^{(1)}\odot M)x+b^{(1)}\right)+b^{(2)}
\]
on the task-specific training dataset $\mathcal{D}^{\mathrm{tr}}_\tau$ using
binary cross-entropy. We define $\varepsilon_\tau$ as the empirical 10th
percentile of the validation losses and retain the successful solutions
satisfying $V_\tau(M)\leq\varepsilon_\tau$.

Each successful solution is represented by its normalized first-layer
importance map
\[
\psi(s_\tau(M))=I_\tau(M)
=
\frac{\left|\widehat{W}^{(1)}_{\tau,M}\odot M\right|}
     {\left\lVert\widehat{W}^{(1)}_{\tau,M}\odot M\right\rVert_\infty},
\]
where $\widehat{W}^{(1)}_{\tau,M}$ denotes the first-layer weights obtained
by training the model with mask $M$ on task $\tau$. Compared with the binary
mask alone, the importance map additionally retains the relative magnitudes
of the learned connections.

We train a CVAE on pairs $(I_\tau(M),\bar g)$, using the normalized scalar
$\bar g=(g-3)/(10-3)$ as the only task condition. The patterns $A$ and $B$
are not provided, since the reference support depends on their relative
displacement but not on their identities.
Given $z\sim p_z=\mathcal{N}(0,\mathbf{I})$, the decoder produces an
importance map $G_\theta(z,\bar g)$. The structure map $\pi$ retains its 96 largest
entries, yielding
\[
\widetilde{M}=\pi(G_\theta(z,\bar g)),
\qquad
q_\theta(\,\cdot\mid \bar g)
=
\bigl(\pi\circ G_\theta(\cdot,\bar g)\bigr)_{\#}p_z.
\]
Thus, every structure sampled from the learned conditional prior has exactly
96 active connections out of 256.


\paragraph*{Gap-held-out meta-test evaluation}
To test whether the generator transfers a structural rule rather than only
reproducing structures observed at particular conditions, we remove complete
gap values from meta-training while keeping the same eight motif pairs in both
splits. Each split contains 48 meta-training and 16 meta-test tasks. We use two
profiles: scalar interpolation, with gaps $5$ and $8$ held out, and one-sided
scalar extrapolation, with gaps $3$ and $4$ held out and gaps $5$--$10$ used
for meta-training.

For each meta-test task, we sample 64 masks from each prior and train fresh
MLP weights on the task-specific training data. We compare the CVAE with an
exact-96 random prior, an unconditional VAE, a train-only conditional mean,
a nearest-wrong-gap CVAE, and the ideal Toeplitz support. For an interpolation
gap with two equally near meta-training neighbours, the wrong-gap result
averages both conditions using the same latent samples as the correct-gap
query. We report BCE, accuracy, and best-permutation IoU. No meta-test labels,
losses, or importance maps are used to construct or select the masks.

\begin{table}[htbp]
  \centering
  \small
  \caption{Gap-held-out meta-test performance. Lower BCE is better; higher
    accuracy and best-permutation IoU are better. Values are means across the
    16 held-out tasks in each profile.}
  \label{tab:gap-ood-meta-test}
  \begin{tabular}{lccc}
    \toprule
    Prior & BCE $\downarrow$ & Accuracy $\uparrow$ & IoU $\uparrow$ \\
    \midrule
    \multicolumn{4}{l}{\emph{Interpolation: held-out gaps $5$ and $8$}} \\
    Random exact-96        & 0.583 & 0.718 & 0.429 \\
    Unconditional VAE      & 0.570 & 0.726 & 0.561 \\
    CVAE, nearest wrong gaps & 0.562 & 0.735 & 0.605 \\
    CVAE                   & 0.561 & 0.736 & 0.608 \\
    Conditional mean       & 0.556 & 0.742 & 0.637 \\
    Ideal Toeplitz support & 0.547 & 0.755 & 1.000 \\
    \midrule
    \multicolumn{4}{l}{\emph{Extrapolation: held-out gaps $3$ and $4$}} \\
    Random exact-96        & 0.582 & 0.713 & 0.437 \\
    Unconditional VAE      & 0.602 & 0.689 & 0.384 \\
    CVAE, nearest wrong gap & 0.587 & 0.709 & 0.418 \\
    CVAE                   & 0.586 & 0.711 & 0.410 \\
    Conditional mean       & 0.577 & 0.722 & 0.525 \\
    Ideal Toeplitz support & 0.566 & 0.740 & 1.000 \\
    \bottomrule
  \end{tabular}
\end{table}

\begin{figure*}[htbp]
  \centering
  \includegraphics[width=\textwidth]{gap_ood_cvae_vs_ideal}
  \caption{Fixed-seed samples from the scalar-conditioned CVAE compared with
    their ideal Toeplitz supports at held-out gaps. The first two columns use
    the interpolation model and the last two use the one-sided extrapolation
    model. Generated masks are aligned to the corresponding ideal masks by an
    optimal permutation of the exchangeable hidden units; IoU is reported
    after alignment. Every displayed mask contains exactly 96 active
    connections.}
  \label{fig:gap-ood-generalization}
\end{figure*}

In the interpolation profile, the CVAE improves over the exact-96 random prior
by $0.018$ accuracy and $0.180$ IoU. The paired accuracy difference has a 95\%
confidence interval of $[0.003,0.033]$ ($p=0.024$). However, the gain is
concentrated at gap $8$: at gap $5$, the CVAE and random prior have comparable
accuracy. The correct-gap query is only marginally better than the symmetric
nearest-wrong-gap aggregate (accuracy $0.736$ versus $0.735$; IoU $0.608$
versus $0.605$), and it remains below the train-only conditional mean. More
stringently, comparing each task with whichever of its equally near wrong-gap
queries has the higher IoU gives a mean margin of $-0.017$; the margin is not
positive for any of the 16 tasks. Thus the apparent aggregate interpolation
gain does not establish reliable use of the scalar condition.

Under one-sided extrapolation, the CVAE does not improve upon the exact-96
random prior: the accuracy difference is $-0.002$ with a 95\% confidence
interval of $[-0.006,0.001]$ ($p=0.191$), and its mean IoU is also lower.
The strict condition-fidelity margin is again negative on average
($-0.009$) and is positive for only 8 of 16 tasks.
Consequently, the experiment shows that the learned generator can produce
useful in-range structure, but it does not yet establish a robust
gap-conditioned inductive bias that extrapolates beyond the meta-training
regime.

\paragraph*{Scope of the structural claim}
The reference support is sufficient, but not uniquely identified by this
restricted task distribution: an audit found alternative supports that solve
some tasks exactly on the admissible inputs. Moreover, resampling this small
input space can repeat strings across data splits. A separate audit of later
task-loss experiments found extensive repetition; overlap for
Table~\ref{tab:gap-ood-meta-test} has not been independently quantified.
We therefore do not use this table to establish generalization to previously
unseen input strings.

\subsection{Pattern-Length Interpolation}
% Numerical source: mds/data/2026-09-08/pattern32_summary.json,
% produced by pattern/length_interp/report.py (5,000 pattern bootstrap draws).

We next test the same prior-extraction approach on linear binary sequences
$x\in\{0,1\}^{32}$. A task is a fixed pattern $p$ of length $k$, with label
$y=1$ exactly when $p$ occurs as a substring. Unlike the motif-pair setting,
zero, one, and multiple occurrences are allowed. A global hash partitions
input strings into training, validation, and test sets independently of the
task. Patterns related by reversal or bit complementation remain in the same
task split. Lengths $3,4,6,8$ are used for prior construction and checkpoint
selection; all 32 patterns of length 5 and all 128 of length 7 are held out.

For each training task, we train 1,000 MLPs with 32 hidden units and random
first-layer masks containing exactly $32k$ connections. The normalized
importance maps of the best 10\%, selected by validation BCE, train a CVAE
conditioned only on the scalar $2(k-3)/5-1$. At evaluation, each of two fitted
CVAE seeds supplies 32 masks per held-out length. Each mask is evaluated with
two fresh MLP initializations, 2,000 Adam steps on 8,192 training examples,
and 2,048 balanced test examples. All methods use the same target cardinality
$32k$; wrong-condition controls also share the latent samples.

\begin{table}[htbp]
  \centering
  \small
  \caption{Pattern-length interpolation: balanced test accuracy (\%).
    Wrong conditions are $k-1$ and $k+1$. The analytic support is followed
    by ordinary MLP training, not analytically assigned weights.}
  \label{tab:pattern-length-interpolation}
  \begin{tabular}{lcc}
    \toprule
    Prior & $k=5$ & $k=7$ \\
    \midrule
    Random exact-$32k$ & 65.95 & 65.90 \\
    CVAE, lower condition & 69.18 & 69.17 \\
    CVAE, upper condition & 69.44 & 69.06 \\
    CVAE, correct condition & 69.34 & 69.16 \\
    Analytic sliding-window support & 73.34 & 71.83 \\
    \bottomrule
  \end{tabular}
\end{table}

The CVAE improves over random masks by 3.39 and 3.26 percentage points at
$k=5$ and $k=7$, respectively. Paired 95\% bootstrap intervals over held-out
patterns are $[2.85,3.95]$ and $[3.00,3.51]$ percentage points, conditional
on the two fitted CVAE seeds. Generated supports show local structure
(Figure~\ref{fig:pattern-length-masks}), yet the correct length condition
performs similarly to either neighbour. Thus this experiment demonstrates
transfer of a useful structural prior to unseen intermediate lengths,
without establishing recovery of the length-dependent support rule.

\begin{figure}[htbp]
  \centering
  \includegraphics[width=0.85\linewidth]{pattern_length_masks}
  \caption{The first evaluated CVAE sample from seed 42 at each held-out
    length, compared with the analytic support. Hidden units are permuted
    only for visual alignment; no sample is selected by its IoU or accuracy.}
  \label{fig:pattern-length-masks}
\end{figure}
